\section{Global shielding and extension to cofaces}\label[appendix]{sec:shield-global}

The local annular cost does not add once for every singular piece:
a concentration argument bounds a new shrinking cost by the maximum
of the old and local form coefficients.  We prove that assertion first,
then transport each chosen addition to cofaces, and finally show that
only finitely many additions are needed on a whole packet.
The finite-label hypotheses in this appendix are extended in
\cref{countable:finite-data,countable:step}.  Each label may already
have its own positive potential floor $p_j^*$; none of the estimates
uses a positive lower bound common to all countably many labels.

We keep the true cuts fixed throughout this section.  Only profiles of
remaining labels, their cores, and their potentials are changed.  In
particular, none of the convergence assertions below concerns minimizers
for a changing pressure.

\begin{definition}[The matching form domain]\label{shield:domain}
Let $E$ be a parent packet of dimension $m$, with true wall $S$ and singular
set $\Sigma$.  A raw piece $Z$ means either $S_{\rm reg}$ or one of the two
joined spaces
\[
 Z=\Omega\mathbin{\cup}_{S_{\rm reg}}
                 (S_{\rm reg}\times[0,\infty)).
\]
The tests on a joined space are compactly supported, piecewise smooth
functions whose two traces agree at the regular seam and whose support
avoids $\Sigma$ and the cylinders over $\Sigma$.  The wall has the
corresponding compact tests away from $\Sigma$.  All measures include the
old warp weight.  Put
\begin{equation}\label{shield:hilbert}
 \mathfrak h_Z(v)=\mathcal N_Z(v)+e_0\norm{v}_{L^2(Z)}^2,
 \qquad e_0=\frac{\lambda_r-\lambda_{r+1}}{100}>0,
\end{equation}
where $\mathcal N$ is the reserved gradient form in \cref{cut:reserve}:
its gradient coefficients are $.8$ on the bulk and $1.7$ on the wall or
cylinder.  Write $\mathcal H_Z$ for the completion of these tests in
$\mathfrak h_Z^{1/2}$.  Pure cylinders $S_{\rm reg}\times\R$ use the
coefficient $1.7$ throughout.  The definitions also apply after adjoining
finitely many Euclidean product coordinates, with the same coefficients
and $e_0$.  Such products are used only to identify limits on translated
windows; they do not append further circle warps or change the weight.
A compact window means a fixed compact ambient base window together with
a bounded cylinder interval, when a cylinder is present.
\end{definition}

The trace matching in this definition is essential.  There is no
Dirichlet condition at the seam.  A pure-cylinder test can be translated
into the cylindrical part of a joined space, with zero bulk part.
Consequently every form inequality on that joined space for potentials
pulled back from its base also holds on the pure cylinder.

\begin{lemma}[Compactness on joined windows]\label{shield:compactness}
Assume $m\ge8$.  On any fixed compact window, a bounded sequence in
$\mathcal H_Z$ has a strongly convergent $L^2$ subsequence.  The following
moving version also holds.  Work in fixed smooth ambient geometry on each expanding window,
as is the case after translation in an exact product end.  Suppose
positive weights and pressures converge smoothly there, and the true
phases and walls converge as in \cref{cut:gmt,cut:end-convergence}.
Consider normalized tests on the corresponding raw pieces $Z_i$, all
supported in a fixed compact window strictly inside a larger window of
convergence.  After passage to a subsequence there is a function $v$ on
the limiting raw piece $Z_\infty$ such that:
\begin{enumerate}[label=\textup{(\roman*)}]
\item On every regular compact subwindow, the tests converge weakly in
$H^1$ and strongly in $L^2$ in the smooth graphical identifications.  At
a limiting regular seam the two traces agree.
\item No $L^2$ mass is lost in neighborhoods of the limiting singular
set.  In particular, integrals of $v_i^2$ converge to those of $v^2$,
including integrals against any fixed bounded continuous ambient
potential and its pullback to the wall and cylinder.
\item The limit belongs to $\mathcal H_{Z_\infty}$ and
\begin{equation}\label{shield:liminf}
 \mathfrak h_{Z_\infty}(v)
       \le\liminf_i\mathfrak h_{Z_i}(v_i).
\end{equation}
\end{enumerate}
These assertions include pure cylinders, translated cylinder intervals,
and Euclidean product limits.  Constants are local to the enlarged
window; no uniform bound among all packets is asserted.
If the limiting wall is smooth on the enlarged window, the same
conclusions follow from ordinary Rellich compactness and the bounded
trace maps in the smooth graphical charts, in every base dimension.
\end{lemma}

\begin{proof}
We give the estimates for tests first; they extend to the completions by
density.  On the enlarged window the weights are uniformly comparable
to a positive constant, the ambient geometry is uniformly smooth, and
the unweighted base mean curvature of the regular walls is bounded.
The latter bound follows from the prescribed weighted mean curvature
and the bounded first derivative of the logarithm of the weight.
The perimeter varifolds have density at least one and no additional
boundary term on this interior window, by \cref{cut:gmt}.  They therefore
meet the hypotheses of the Sobolev inequality used below.
No bound on the individual principal curvatures is needed.

We use the submanifold Sobolev inequality in the following precise form.
For a $d$-dimensional Euclidean rectifiable varifold of density at least
one with generalized mean curvature $H$, and a nonnegative compactly
supported test $f$,
\begin{equation}\label{shield:submanifold-sobolev}
 \left(\int f^{d/(d-1)}\right)^{(d-1)/d}
 \le C_d\int\bigl(\abs{\nabla f}+\abs H f\bigr).
\end{equation}
This holds in arbitrary codimension; see \citet[Chapter~4, Section~5,
Theorem~5.7]{simonGMT}. This is the submanifold Sobolev inequality
of Michael and Simon \citep{michaelSimon1973}.  A smooth extension of the fixed ambient chart
metric to a compact manifold and a fixed smooth isometric embedding
\citep{nash1956} give the corresponding local inequality in our charts.
The extra mean curvature of the embedding is bounded on this fixed
compact geometry.  For the moving-domain assertion the embedding is chosen once in the
fixed ambient geometry; its constants are therefore independent of the
wall in the sequence.  Applying
\cref{shield:submanifold-sobolev} to a power of the test gives the usual
$H^1$ Sobolev exponent on the wall.  Applied to
$S_{\rm reg}\times\R$, it gives, with
\[
 q=\frac{2m}{m-2},\qquad b=\frac{2(m-1)}{m-2},
 \qquad 2(b-1)=q,
\]
a local $L^q$ bound on the cylinder.  A test on the half-cylinder is
reflected evenly at its base before this application.  The wall alone
has the stronger exponent $2(m-1)/(m-3)>2$.

Here and below cutoffs are taken on the enlarged window and are pulled
back to both seam sides.  The cutoff errors are bounded by the local
$L^2$ norm.  Normalize $\mathfrak h_Z(v)\le1$.  Choose a cylinder cutoff
which is one at $s=0$ and vanishes at the outer end of a fixed bounded
interval.  The fundamental theorem of calculus and H\"older give
\begin{equation}\label{shield:trace}
 \int_{S_{\rm reg}}\abs{v(0)}^b
 \le C\int_{S\times[0,L]}
       \bigl(\abs v^b+\abs v^{b-1}\abs{\partial_s v}\bigr)
 \le C.
\end{equation}
The last step uses the cylinder $L^q$ bound and $2(b-1)=q$.

The bulk need not have a uniform smooth boundary, so we use a BV
calculation there.  Extend $\abs{v_{\rm bulk}}^b$ by zero in the ambient
chart, including its compact cutoff.  For an initial test its support
misses the singular set, and the BV jump formula gives precisely the
boundary contribution $\int_{S_{\rm reg}}\abs{v(0)}^b$ in addition to the
interior variation.  This is also obtained directly by integration by
parts against unit vector fields on the regular side.  The ambient
$L^1$ Sobolev inequality and Cauchy--Schwarz therefore yield, for
$J=\int_\Omega\abs v^q$,
\begin{equation}\label{shield:bulk-bv}
 J^{(m-1)/m}
 \le C\left(
       \int_\Omega\abs v^b+
       \int_\Omega\abs v^{b-1}\abs{\nabla v}
       +\int_{S_{\rm reg}}\abs{v(0)}^b\right)
 \le C(1+J^{1/2}).
\end{equation}
Here $bm/(m-1)=q$; the first integral on the right is at most
$J^{1/2}\norm v_2$, because $b=(q+2)/2$.  Since $(m-1)/m>1/2$, \cref{shield:bulk-bv} bounds $J$.
Rescaling removes the normalization.  Thus all pieces have a local
$L^p$ bound for some $p>2$ controlled by their matching form norm.
This argument extends a power as a BV function; it constructs no
$H^1$ extension operator for the bulk.

We next exclude concentration in the moving version.  On a slightly
larger compact window let $F$ be the compact portion of the limiting
singular set.  Its ambient volume and its limiting perimeter measure
are zero.  For a closed $\delta$-neighborhood of $F$, perimeter measure
convergence and the closed-set upper bound for weakly convergent
measures imply
\[
 \limsup_i\mu_{S_i}(\overline{U_\delta(F)})
       \le \mu_{S_\infty}(\overline{U_\delta(F)}),
\]
with a fixed outer compact cutoff if necessary.  The right side tends
to zero as $\delta\downarrow0$.  The same conclusion for bulk measure
follows from ambient volume bounds, and for a cylinder follows by
multiplying the wall measure by the bounded interval length.  H\"older
and the local $L^p$ estimates give
\[
 \int_{A_i}\abs{v_i}^2
 \le C\mu_i(A_i)^{1-2/p}.
\]
Thus the order is first $i\to\infty$ at a fixed singular neighborhood,
then $\delta\downarrow0$.  Neither disappearing components nor a thin
boundary layer can retain $L^2$ mass.

Off $F$, graphical convergence and the absence of stray wall support
supply smooth charts for both bulk phases and their seams.  Rellich
compactness in those charts gives the asserted weak $H^1$ and strong
$L^2$ limits.  The trace map on each fixed regular chart is bounded and
linear; weak convergence of both traces preserves their equality.
Combining these regular limits with the no-concentration conclusion
proves all the stated $L^2$ and bounded-potential convergence.  Exhausting
the regular pieces gives lower semicontinuity of their gradient
integrals.

It remains to justify membership in the limit completion, rather than
assume it.  On this one limiting geometry, monotonicity and the local
area bound give
\[
 \mu_S(B_{2a}(x))\le C a^{m-1}
\]
for sufficiently small balls in the fixed window.  Since
$\dim_{\rm H}\Sigma\le m-8<m-3$, a compact singular portion has ball
covers, inside arbitrarily small neighborhoods, with
$\sum_i a_i^{m-3}$ arbitrarily small.  For cutoffs equal to one on these
balls, supported on their doubles, and with gradients at most $C/a_i$,
the maximum of the cutoffs has squared-gradient integral at most
\[
 C\sum_i a_i^{m-3}
\]
on the wall and on bounded cylinder intervals.  The bulk cost is at
most $C\sum_i a_i^{m-2}$ and is smaller.  A finite subcover suffices.
Use the same base cutoff on the bulk and cylinder, preserving seam
matching.  First truncate the limit function to a bounded function;
then multiply it by one minus these cutoffs.  The derivative of the
cutoff has vanishing energy after truncation, and the other error tends
to zero by absolute continuity of the already finite gradient
integrals.  Its $L^2$ error tends to zero as well.  The resulting functions
miss the singular set.  Matching traces glue their two $H^1$ pieces
across every regular seam chart; ordinary mollification in such charts,
and a partition of unity, approximate them by the defining tests.
This proves membership and \cref{shield:liminf}.  Capacity is used only
on the limit; no cheap cutoff on every approximating singular set, and
no convergence theorem for the entire moving Sobolev spaces, is needed.
The same proof applies to pure cylinders and product limits.
For a product with $a$ extra coordinates, use ambient dimension $m+a$
and singular-dimension bound $m+a-8$ in the preceding exponents;
the strict codimension needed for the capacity estimate is unchanged.
Thus this argument supplies exactly regular graphical identifications,
strong $L^2$ convergence without singular concentration, matching seam
traces, and the lower semicontinuity used in the later form transfers.
It does not assert smooth convergence at singular points.
\end{proof}

\begin{lemma}[Maximum of successive form costs]\label{shield:maximum}
Fix a packet and its true pieces.  Let $V\ge0$ be a fixed bounded
potential, pulled back to cylinders, satisfying
\begin{equation}\label{shield:old-bound}
 \int_Z Vv^2\le\theta\mathfrak h_Z(v)
\end{equation}
on its wall and both joined sides.  Let $F\subset\Sigma$ be a compact
piece and let $A_\eps=1.16D_\eps$ be one of the nonnegative shell
potentials of \cref{local:shell,local:sharp}.  In particular its base
support lies in a fixed compact set and shrinks to $F$, and
\begin{equation}\label{shield:local-input}
 \int_Z A_\eps v^2\le .78\mathcal N_Z(v)+C_F\norm v_2^2
\end{equation}
for all sufficiently small $\eps$, on all three pieces.
For every $\eta>0$, a sufficiently small $\eps>0$ gives simultaneously
\begin{equation}\label{shield:maximum-bound}
 \int_Z(V+A_\eps)v^2
 \le\bigl(\max\{\theta,.78\}+\eta\bigr)\mathfrak h_Z(v).
\end{equation}
It suffices that $V$ be bounded on the compact base windows used below.
In particular the assertion applies to a finite sum of compact
additions on the initial radial parent.
\end{lemma}

\begin{proof}
First $\int A_\eps v^2\to0$ for every fixed $v\in\mathcal H_Z$.
Indeed \cref{shield:local-input} bounds these forms uniformly by a
constant times $\mathfrak h_Z$, whereas for a defining test the integral
is zero for small $\eps$, since its support misses $F$.  Density and
Cauchy--Schwarz for the nonnegative form give the assertion.

We detail the localization needed on noncompact pieces.  If
$\sum_\alpha\chi_\alpha^2=1$ is a locally finite Lipschitz partition
which has the same seam traces, then
\begin{equation}\label{shield:ims}
 \sum_\alpha\mathfrak h_Z(\chi_\alpha v)
 =\mathfrak h_Z(v)+
   \int_Z c_Z\sum_\alpha\abs{\nabla\chi_\alpha}^2v^2,
 \qquad
 c_Z=\begin{cases}.8&\text{on bulk},\\1.7&\text{on wall or cylinder}.
 \end{cases}
\end{equation}
The mass term $e_0\norm v_2^2$ allows the last integral to be at most
$\tau\mathfrak h_Z(v)$ for any prescribed $\tau>0$.  For example use a
sine/cosine pair of a Lipschitz cutoff angle, equal to $\pi/2$ on a
fixed compact base set containing $\supp A_\eps$ for all small $\eps$,
and zero outside a sufficiently large compact neighborhood.  Its transition width makes its derivative
arbitrarily small.  The cosine part sees only $V$. Its bound is the assumed global
form inequality for $V$; no compact support assumption on $V$ is needed.
The sine part has compact base support, where the old potential is bounded.
On a joined side
also use a squared partition into cylinder intervals of a sufficiently
large fixed length.  The first member is identically one near $s=0$
and extends as one on the bulk; all the others vanish near the seam
and extend as zero there.  Their squared gradients have an arbitrarily
small sum.  All cutoffs in the base are pulled back to the cylinder.

Choose the partition error so small that a violation of
\cref{shield:maximum-bound} forces one localized test to have quotient
larger than $\max\{\theta,.78\}+\eta/2$.  This follows directly by
summing the potential integrals and \cref{shield:ims}; for instance it
suffices to have
$(\max\{\theta,.78\}+\eta/2)(1+\tau)
 <\max\{\theta,.78\}+\eta$.
The selected test has compact base support and a bounded cylinder
interval.  Passing to a subsequence, its interval either stays in a
fixed bounded range, or escapes to infinity.  In the latter case it is
translated to a fixed interval on the pure cylinder.  Both the old
bound and \cref{shield:local-input} hold there by translating its tests
into the original joined side.  Thus in all cases the selected sequence
lies in one fixed Hilbert space, in a fixed compact window.

Suppose now that $\eps_i\downarrow0$ gives such violating tests $v_i$,
normalized by $\mathfrak h(v_i)=1$.  By \cref{shield:compactness} and
weak Hilbert compactness, $v_i\rightharpoonup v$ in this fixed space and
$v_i\to v$ strongly in $L^2$.  Write $w_i=v_i-v$.  Then
\[
 \mathfrak h(w_i)=1-\mathfrak h(v)+o(1),\qquad \norm{w_i}_2\to0.
\]
The bounded old potential satisfies
$\int Vv_i^2\to\int Vv^2\le\theta\mathfrak h(v)$.  The new form
vanishes on $v$, and its mixed term vanishes by Cauchy--Schwarz and the
uniform bound on $w_i$.  Consequently \cref{shield:local-input} gives
\begin{align*}
 \limsup_i\int(V+A_{\eps_i})v_i^2
 &\le\theta\mathfrak h(v)+.78\limsup_i\mathcal N(w_i)\\
 &\le\theta\mathfrak h(v)+.78(1-\mathfrak h(v))
 \le\max\{\theta,.78\}.
\end{align*}
This contradicts the selected quotient.  There are only three raw
pieces, so the same sufficiently small scale works for all of them.
\end{proof}

We next extend a single chosen addition.  The scale in
\cref{shield:maximum} is fixed before this next limit is taken.

\begin{lemma}[Extension with late fronts]\label{shield:extension}
Let $B$ be a packet, let $D\ge0$ be a fixed smooth compactly supported
potential addition on $B$, and let $V$ denote the old cumulative
potential on every packet.  Suppose these old potentials are exactly
compatible and obey the common form bound $\theta\mathfrak h$.
Suppose also that
\[
 \int(V_B+1.16D)v^2\le\beta\mathfrak h_B(v),\qquad \beta\ge\theta,
\]
on the wall and joined sides over $B$.  In each of the finitely many
coface occurrences $B\times(T,\infty)^d\subset E$, introduce a smooth
potential $D_E$ with
\begin{equation}\label{shield:extension-support}
 0\le D_E\le D\circ\pi_B,\qquad
 \supp D_E\subset K\times\prod_{i=1}^d[R_i,\infty),
 \qquad K\Subset B.
\end{equation}
Use consistent functions on the normal coordinates identified by
commuting face chains, so that these extensions are compatible.
For every $\eta>0$ all starts can be chosen sufficiently large that
\begin{equation}\label{shield:extended-bound}
 \int(V_E+1.16D_E)v^2\le(\beta+\eta)\mathfrak h_E(v)
\end{equation}
on every affected raw piece.  The conclusion concerns the same fixed
true cuts; it uses their product convergence from
\cref{cut:end-convergence}.
\end{lemma}

\begin{proof}
Only finitely many occurrences are involved by \cref{packet:category}.
Here and below a source is extended along the directed face-to-coface
relations, through their finite depth.  Closing under arbitrary
alternating face and coface walks is neither required nor performed.
For a common normal coordinate take the largest of the finitely many
requested starts, so commuting chains give identical occurrences.
It is enough to prove the assertion for any sequence of their fronts
with $\min_iR_i\to\infty$, passing to a fixed occurrence and piece if
needed.  Choose every localization width below using $\eta$ and $e_0$
first, independently of the front positions.

Separate the new support from its complement with a sine/cosine pair
as in \cref{shield:ims}.  More explicitly, an angle can be $\pi/2$
times a product of a compact base cutoff, equal to one on $K$, and
normal cutoffs which equal one on $[R_i,\infty)$ and vanish before
$R_i-L$.  Enlarge the compact base window and $L$ until the sum of
squared derivatives is sufficiently small.  Extend the angle by zero
outside the product sector.  Taking the fronts beyond all collar
thresholds makes this extension well defined.  The cosine test misses
the addition and has the old bound $\theta$.  Partition the sine part
further into fixed large boxes in every normal coordinate, and into
fixed large intervals in cylinder height when applicable.  The total
localization error is as small a fraction of $\mathfrak h$ as required.

If \cref{shield:extended-bound} failed along fronts tending out, the
same quotient selection as in \cref{shield:maximum} would give
normalized tests on selected boxes with quotient at least
$\beta+\eta/2$.  Such a box must meet the new potential support,
since otherwise the old bound $\theta\le\beta$ applies.  Its center
therefore has
\[
 t_i\ge R_i-O(L)\longrightarrow\infty
 \quad\text{for every }i=1,\ldots,d.
\]
Its base window is fixed and compact.  Recenter every normal coordinate.
A bounded cylinder interval stays in a bounded-height joined window,
or, after another subsequence, is translated to a pure cylinder.  The
ambient geometry on the recentered windows is exactly the product
geometry, and the true cuts converge to those over $B$ times $\R^d$.
The old potential is exactly $V_B\circ\pi_B$ there, because the fronts
also exceed its fixed product thresholds.  The new potential is bounded
above by the fixed bounded function $D\circ\pi_B$.

Apply \cref{shield:compactness}.  For its limit $v$,
\begin{equation}\label{shield:extension-limit}
 \limsup\int(V_E+1.16D_E)v_E^2
 \le\int(V_B+1.16D)\circ\pi_B\,v^2,
 \qquad \mathfrak h_{\rm lim}(v)\le1.
\end{equation}
The face form bound applies to the right side by slicing the $\R^d$
coordinates.  For defining tests this is Fubini's theorem: the gradient
terms in the extra directions are nonnegative, and all coefficients
$.8,1.7,e_0$ are the same in every base dimension.  Approximation in the
limit completion, proved in \cref{shield:compactness}, gives the result
for $v$.  In a pure-cylinder limit use the face joined bound after
translation in cylinder height.  Thus the right side of
\cref{shield:extension-limit} is at most $\beta$, a contradiction.
Finite incidence permits simultaneous choices for all occurrences.
This limit compares bounded potential integrals on moving domains;
it uses no subtraction of vectors belonging to different form spaces.
\end{proof}

\begin{lemma}[Compatible caps and cores]\label{shield:caps}
A compact forcing update from \cref{local:shell} extends to all its
cofaces so as to satisfy \cref{shield:extension-support}, preserve the
profile gradient inequality with coefficient $a_r$, and enlarge only
the designated cores.  The extension is an exact product sufficiently
far in every old face collar and is made once per occurrence, independent
of the order of commuting face relations.
\end{lemma}

\begin{proof}
Consider a plus addition on $B$, with forcing profile $u$ and shell
potential $D$.  Choose a compact base window $K\Subset\{h_j>t_j\}$
containing both $\overline{\{u> -\infty\}}$ and $\supp D$ in its
interior.  The first set includes the inner pole region; keeping only
$\supp D$, which is annular, would not suffice for this extension.
In a coface product over $K$ put
\begin{equation}\label{shield:cap-formula}
 u_E(x,t)=\min\{u(x),w_1(t_1),\ldots,w_d(t_d)\},\qquad
 D_E(x,t)=D(x)\prod_{i=1}^d\chi_i(t_i).
\end{equation}
Each cap $w_i$ is $-\infty$ up to a late front, then crosses to
$+\infty$, and remains there, with
\[
 2\abs{w_i'}\le .58w_i^2+p_j^*
\]
on its finite band.  For example integrate $2w_i'=.58w_i^2+p_j^*$ on the finite
band and extend by its poles.  The crossing length is
$2\pi/\sqrt{.58p_j^*}<\infty$; finite truncations remain locally
Lipschitz at both ends.  Only this label's positive floor is used.
Each smooth cutoff $\chi_i$ takes values in
$[0,1]$ and \emph{finishes} increasing to one before $w_i$ begins its
crossing.  Choose these cutoff fronts to satisfy
\cref{shield:extension} for the already fixed $D$ and shell scale.
Outside this product occurrence use the inactive value $-\infty$.
The old profile is replaced by its maximum with $u_E$.

For locally Lipschitz finite truncations, a finite minimum or maximum
has the derivative of one active branch almost everywhere (on a tie,
the branch derivatives agree almost everywhere on the tie set).
If the active finite branch is $u$, every cap has already left its
$-\infty$ plateau; hence every $\chi_i=1$ and the entire shell
potential $D$ is available.  If it is a cap, its derivative cost is
paid by the preexisting baseline:
\[
 .58w_i^2+p_j^*\le a_rw_i^2+\Pi_j+D_E.
\]
The old branch retains its old inequality.  Thus no sum of normal
cap derivatives, and no extra cap potential, enters the budget.
For a minus update negate $u_E$ and take a minimum with the old profile.
The strict anchor regions for opposite signs are disjoint and contain
no opposite old poles by \cref{local:anchors}.

Write $T_i^+$ for the end of the $i$th cap crossing.  Enlarge the
plus core by the open set
\[
 \{x:\delta_B(x)<\varepsilon/2\}\times
       \prod_i(T_i^++b_i,\infty),\qquad b_i>0.
\]
The properly embedded closed tails ensure that its closure in the
coface has $\delta_B\le\varepsilon/2$ and $t_i\ge T_i^++b_i$.
It is therefore contained in the open set
$\{\delta_B<\varepsilon\}\times\prod_i(T_i^+,\infty)$, where every
branch of the minimum is $+\infty$.  This proves closure confinement,
including at the cap fronts, with the corresponding statement for a
minus core.  The front positions are also chosen beyond the exact anchor-product
thresholds, so its support retains the designated strict anchor sign.

Use the same caps and cutoffs for the same normal coordinates in every
identified occurrence.  On a face obtained by hitting one of these
coordinates, that coordinate's cap is eventually $+\infty$ and its
cutoff is one, and that factor disappears from
\cref{shield:cap-formula}.  If an unrelated coordinate is hit, an
intersection sufficiently far in its collar would require $x$ to
escape into a further face of $B$, which the compact support forbids.
Finite incidence allows all relevant intersection thresholds to be
exceeded at once.  These observations prove compatibility for all
chains and show why a commuting product occurrence is not counted
more than once.
\end{proof}

\begin{proposition}[The global shielding output]\label{shield:global}
At a fixed stage $r$, let $j_0$ be the current label and keep its true
minimizers fixed.  There are updates of every remaining profile and
core, and compatible smooth additions $D_j\ge0$, with the following
properties on every parent packet $E$ of dimension $m$:
\begin{enumerate}[label=\textup{(\roman*)}]
\item The updated profiles have locally Lipschitz finite truncations
and satisfy, almost everywhere on their finite bands,
\begin{equation}\label{shield:profile-output}
 2\abs{\nabla\mu'_j}\le a_r(\mu'_j)^2+\Pi_j+D_j.
\end{equation}
All these data are exact compatible products.  Each core only enlarges,
and its closure remains in the interior of its designated updated pole
set.  The original forced sets and the anchor sign restrictions are
preserved.
\item With $G'_j=E\setminus(U_j^{\prime+}\cup U_j^{\prime-})$, the set
\begin{equation}\label{shield:guard-output}
 H_E^*=\left\{x:\#\{j\ne j_0:x\in G'_j\}\ge m-2\right\}
\end{equation}
is closed and misses $\Sigma_E$.  In an exact collar over a face $B$,
its projection is contained in $H_B^*$.  This holds for every further
face and every simultaneous product limit.
\item For one common $\theta_*<.9$, all walls and joined sides satisfy
\begin{equation}\label{shield:form-output}
 \int 1.16\sum_{j\ne j_0}D_jv^2
 \le\theta_*\bigl(\mathcal N(v)+e_0\norm v_2^2\bigr)
 \le .9\bigl(\mathcal N(v)+e_0\norm v_2^2\bigr).
\end{equation}
The additions are bounded with their derivatives on each post-radial
packet.  There are only finitely many local additions on each such
packet and only finitely many extensions into it.
\end{enumerate}
For the initial radial parent the same conclusions hold with locally
finite additions over its compact transition shells; boundedness is
asserted on the corresponding compact blocks and
\cref{shield:form-output} holds on every compact test.  No convergence
assertion at an end of the original manifold is needed.
\end{proposition}

\begin{proof}
First consider the packet system after the radial cut.  Process its
finite set of depths from the deepest faces upward.  The anchors and
thresholds of \cref{local:anchors} are fixed once, before any updates.
Initially the cumulative addition is zero.  Maintain a single common
upper form bound at least $.78$; it is harmless to start this bookkeeping
bound at $.78$ although the initial form vanishes.  Enumerate the
countable pairs consisting of a packet and a possible successive local
addition number by distinct integers $\nu\ge1$.  Reserve
$\gamma_\nu=.1\,2^{-\nu}$ for that pair, allocating at most half to
\cref{shield:maximum} and half to the simultaneous extensions of
\cref{shield:extension}.  The common bound after any collection of
completed updates is at most
\[
 .78+\sum_{\nu\text{ used}}\gamma_\nu\le .88<.9.
\]
We may consequently take $\theta_*=.88$.  There is no requirement
that a shell scale or its extension fronts be uniform in the packet
or in the addition number.

Suppose all strict further faces of $E$ have been treated and their
updates extended.  The remaining risky set
\[
 K_E=\Sigma_E\cap H_E^*
\]
is compact.  To prove this, if points in it escaped, the packet escape
property \cref{packet:escape} would give a subsequence over a deeper
face $B$, with compact convergent base coordinates and every escaping
normal coordinate tending to infinity.  The updated cores are exact
products there.  There are finitely many relevant labels, so pass to a further
subsequence with a fixed subset of $m-2$ undecided labels.
Membership in each of their closed sets $G'_j$ passes to the base
limit.  Thus the limit base point has at least $m-2$ undecided labels.
This implies membership in $H_B^*$, because
$\dim B-2\le m-2$.  By \cref{cut:end-convergence} the translated true
cuts converge to the true cut over $B$ times the normal Euclidean
factor.  Support convergence puts the limiting base point on the true
boundary over $B$.  A regular limiting boundary point would have a neighborhood
of smooth graphical convergence, which contains no singular points of
the approximating walls.  The limit must therefore be singular on the
face.  This contradicts its completed shielding.  Thus $K_E$ has no
escaping sequence; it is closed, and the compact-part property proves
compactness.  If $m<8$ it is already empty.

Denote this initial compact set by $K_E^{(0)}$.  For each remaining
label and each strict anchor sign, cover $K_E^{(0)}$ in that open
region by the countable compact flat pieces of \cref{local:flat-pieces}.
These pieces need not themselves be an open cover.  Enumerate the
signed pieces; their eventual inner forcing neighborhoods will provide
the open cover needed for termination.  For each,
choose its shell sufficiently small by \cref{shield:maximum}, then
hold that scale fixed and extend its caps, potential, and core by
\cref{shield:extension,shield:caps}.  Supports can be as small as needed
to stay within the prescribed strict anchor region.  Opposite updates
for one label therefore never conflict.  Profiles retain
\cref{shield:profile-output} by their active-branch inequalities, and
cores enlarge monotonically.

To prove finite termination, after $N$ additions set
\[
 O_N=\{x\in E:\#\{j\ne j_0:x\in G_j^{(N)}\}<m-2\},
 \qquad K_E^{(N)}=K_E^{(0)}\setminus O_N.
\]
Each $O_N$ is open: with finitely many labels its complement is a
finite union of intersections of closed $G_j^{(N)}$.  Core enlargement
implies $O_N\subset O_{N+1}$.  Fix $x\in K_E^{(0)}$.  For each of its
strict anchors, some enumerated compact piece contains $x$, and the
addition on that piece puts $x$ in an open enlarged core, irrespective
of how small its positive shell scale was chosen.  Only finitely many
anchors are involved.  Once those additions have occurred, only
threshold-equality labels can remain undecided at $x$.  By
\cref{local:anchors} there are at most $m-8<m-2$ of those.
Thus $x\in O_N$ for some finite $N$.  The increasing open sets $O_N$
cover $K_E^{(0)}$, so compactness supplies a finite subcover; taking
its largest index gives $K_E^{(N)}=\varnothing$.  Stop there.
This argument proves an actual finite stopping time before forming any
infinite sum of profiles.  All subsequent updates only shrink the guard.

There may be countably many packets at one depth.  An originating packet
has finitely many directed coface occurrences, so each individual
extension has already been chosen in finitely many places.  Conversely,
a fixed target packet has finitely many source faces, including through
chains, so only finitely many packets at that depth affect it.  Each of
those sources required finitely many additions.  This uses both
directions of finite incidence from \cref{packet:category}.
Its updated functions are finite sums or finite minima and maxima;
there is no infinite limiting profile on that packet.  Its product
thresholds are bounded by finitely many chosen thresholds, and all its
potential derivatives remain bounded.  The common form bound after a
whole depth follows from the same summable allowance.  This justifies
proceeding to the next depth, and proves the asserted smoothness,
compatibility, and form bounds throughout the system.

The new guard is closed because the number of labels is finite and
each $G'_j$ is closed.  On a far collar, exact product cores give the
same count upstairs and downstairs.  Lower face dimension weakens the
required count, proving the projection assertion in (ii).  These are
precisely the closed-set and product-limit properties needed for the
guarded regularization in \cref{sec:regularization}; a quantitative
smooth tube is obtained there from regular graphical convergence.

For the initial radial cut, process its locally finite compact
transition shells in order, choosing the supports of all additions
inside buffered neighborhoods of their respective shells.
By \cref{packet:initialization,cut:barriers}, the barrier subbands
separate these neighborhoods, and the true boundary lies in
their interiors.  The true singular set in a shell is closed inside its compact buffered
annulus, and the guard is closed.  Their intersection is therefore
compact.  Apply the same increasing-open-cover argument there, keeping
the closure of every forcing region inside this one shell buffer.
There is no coface extension at this step.  Before any addition only
finitely many previous additions have
been made, so their total is a bounded compactly supported potential
and \cref{shield:maximum} applies.  The base localization is obtained
from a proper distance cutoff; all estimates after localization use
only fixed compact data.  Use summable allowances as before.  Local
finiteness leaves smooth potentials and well-defined profiles, and a
compact test sees only finitely many shell additions.  The form bound
therefore persists after all shells are treated.  This proves the
radial assertion and completes the construction.
\end{proof}
